\chapter{Capturing and Monitoring the Network}\label{ch:nw:capturing-and-monitoring-the-network}

In this chapter's simulated cage the strategy's own log adds up, at the median, to \qty{1.5}{\micro\second} between decoding a
packet and handing the order it caused to the network card. The capture taken at the cage's taps says that the order's frame left
the order handoff \qty{4.4}{\micro\second} after the triggering packet's first byte reached the market-data handoff. The difference is
the part of the path the program cannot see: the fibre and the layer-1 devices, the card's receive path into memory, and its
transmit path out. At the 99th percentile the capture says \qty{38}{\micro\second}, and no log written by the program could have
said which part of the path owns it.

This chapter is about the firm watching its own traffic from the outside. It copies every frame at the handoffs (chapter~2's taps),
timestamps it against the clocks of chapter~4, stores it, and answers from the capture alone the questions the firm asks every
day: how long did we take, what did we lose, when did the port fill.

\section{Taps, capture appliances and capture files}

\begin{definition}[Packet capture, capture appliance]\label{def:nw:capturing-and-monitoring-the-network:capture}
A \emph{packet capture}\index{packet capture} is a record of the frames that crossed a point of a network, each with the time it
crossed. A \emph{capture appliance}\index{capture appliance} is a dedicated device that receives copies of traffic from taps,
timestamps each frame against a disciplined clock, and writes the frames to storage at the \omterm{def:nw:networking-for-trading:ser}{line rate} of its inputs without losing
any.
\end{definition}

\begin{definition}[Packet capture file, timestamp trailer]\label{def:nw:capturing-and-monitoring-the-network:trailer}
A \emph{packet capture file}\index{packet capture file} stores a capture as a sequence of records, each a timestamp, a length and the
frame's bytes; the libpcap format and its successor pcapng are the common ones, with nanosecond timestamps in their newer variants.
A \emph{timestamp trailer}\index{timestamp trailer} is a block of bytes that a tapping or switching device appends to the end of a
copied frame, carrying the time at which the device saw the frame and where it saw it (device and port), so that the timestamp
travels with the frame to wherever it is stored.
\end{definition}

The timestamp that matters is the one taken where the frame crossed the tap, not the one the storage server writes when the copy
reaches it: between the two lie an aggregation switch, a queue and the server's own card. A trailer fixes that. It is taken on the
wire, in the frame's first bytes, and it rides the copy through any later aggregation; the storage server keeps it and can ignore its
own time. \Cref{dat:nw:capturing-and-monitoring-the-network:trailers} describes one vendor's published layout; the book's
\texttt{firm.wirecap} uses a 16-byte trailer of its own in the same spirit.

\begin{omfigure}
\begin{tikzpicture}[font=\footnotesize,b/.style={draw,rounded corners=2pt,minimum height=0.6cm,align=center,inner sep=3pt},>=Stealth]
\node[b,fill=gray!10] (A) at (0,1.6) {handoff A};
\node[b,fill=gray!10] (B) at (0,0.5) {handoff B};
\node[b,fill=gray!10] (O) at (0,-0.6) {order handoff};
\foreach \n/\y in {tA/1.6,tB/0.5,tO/-0.6} {\node[b,fill=omExo!15] (\n) at (2.5,\y) {tap};}
\node[b,fill=omDef!15,minimum height=1.2cm] (cage) at (6.2,0.5) {cage devices\\ and servers};
\node[b,fill=omMeth!20,minimum height=1.0cm,minimum width=2.4cm] (agg) at (3.9,-2.6) {timestamping\\ aggregation};
\node[b,fill=omThm!12,minimum height=1.0cm] (app) at (9.6,-2.6) {capture appliance:\\ storage, analysis};
\node[b,fill=gray!8] (gm) at (9.6,-0.9) {grandmaster\\ (chapter~4)};
\foreach \a/\b in {A/tA,B/tB,O/tO} {\draw[thick] (\a) -- (\b);}
\foreach \a in {tA,tB,tO} {\draw[thick] (\a.east) -- (cage.west);}
\foreach \a/\x in {tA/3.3,tB/3.55,tO/3.8} {\draw[->,gray,dashed] (\a.south) -- ++(0,-0.25) -| (\x,-2.1);}
\draw[->,thick] (agg) -- node[above,font=\scriptsize]{copies with trailers} (app);
\draw[->,gray] (gm) -- node[right,font=\scriptsize]{time} (app);
\draw[->,gray] (gm.west) -| (4.8,-2.1);
\end{tikzpicture}

\omcaption{Watching a cage from outside. Taps on every handoff copy both directions; a timestamping device appends a trailer to each copy
(time, device, port) against the grandmaster's time and aggregates the copies; the \omterm{def:nw:capturing-and-monitoring-the-network:capture}{capture appliance} stores and analyses them. No
program in the cage is involved.}\label{fig:nw:capturing-and-monitoring-the-network:taps}
\end{omfigure}

\begin{omfigure}
\begin{tikzpicture}[font=\scriptsize,f/.style={draw,minimum height=0.8cm,align=center,inner sep=1pt}]
\node[f,fill=omDef!12,minimum width=4.2cm] (fr) at (0,0) {original frame (as on the wire)};
\node[f,fill=omMeth!20,minimum width=1.4cm,anchor=west] (d) at (fr.east) {device\\ 2 bytes};
\node[f,fill=omMeth!20,minimum width=1.4cm,anchor=west] (p) at (d.east) {port\\ 2 bytes};
\node[f,fill=omMeth!20,minimum width=1.4cm,anchor=west] (fl) at (p.east) {flags\\ 4 bytes};
\node[f,fill=omThm!20,minimum width=1.5cm,anchor=west] (s) at (fl.east) {seconds\\ 4 bytes};
\node[f,fill=omThm!20,minimum width=1.7cm,anchor=west] (n) at (s.east) {nanoseconds\\ 4 bytes};
\draw[decorate,decoration={brace,amplitude=4pt,mirror}] (d.south west) -- (n.south east)
  node[midway,below=5pt] {16-byte trailer appended by the tapping device (the book's layout)};
\end{tikzpicture}

\omcaption{The \omterm{def:nw:capturing-and-monitoring-the-network:trailer}{timestamp trailer} of \texttt{firm.wirecap}: the device and port that saw the frame and the full UTC time in seconds and
nanoseconds, appended after the frame. Vendors' trailers carry the same information in their own layouts.}\label{fig:nw:capturing-and-monitoring-the-network:trailer}
\end{omfigure}

\begin{dated}{2026-09}{A published trailer}\label{dat:nw:capturing-and-monitoring-the-network:trailers}
Arista documents that its 7130 devices insert a self-contained 16-byte timestamp header at the end of each copied \omterm{def:nw:networking-for-trading:frame}{Ethernet frame},
recompute its check sequence, and carry 64 bits of seconds and nanoseconds, so that ``full UTC time'' needs no keyframes; the header
names the device and port, and a flag says whether the original frame's check sequence was valid. Timestamps are taken ``on the first
byte of a packet on the ingress port'', with the front-panel delay removed, and the document gives $\pm$75~picoseconds of accuracy
when the device is synchronised by a pulse per second. Other devices overwrite a field of the frame with a shorter timestamp and send
periodic keyframes to relate it to UTC.
\end{dated}

\omcode{code/firm/wirecap/firm_wirecap.py}{75}{92}{Writing and reading a nanosecond pcap file: a 24-byte file header, then per frame a
16-byte record header and the bytes.}\label{lst:nw:capturing-and-monitoring-the-network:pcap}

\section{Measuring wire-to-wire latency}

Wire-to-wire latency (One Quant Book~13, chapter~1) is measured on the wire: from the first byte of a market-data frame at the firm's
handoff to the first byte of the order it caused at the order handoff. From a capture it is a subtraction, once each order has been
paired with its cause.

\begin{definition}[Trigger packet]\label{def:nw:capturing-and-monitoring-the-network:trigger}
The \emph{trigger packet}\index{trigger packet} of an order is the market-data packet whose content caused the firm to send it. Its
capture time is the start of the order's wire-to-wire latency.
\end{definition}

The pairing is the hard part. A capture sees two streams of packets, and nothing on the wire says which packet caused which order.
The naive rule, ``the latest packet with a triggering message before the order'', works when triggers are rare and the firm is fast;
it fails exactly when latency is worth measuring, in bursts, when several triggers precede an order by less than the firm's own latency.
The reliable rule is to carry the cause in the order: the firm puts the \omterm{def:nw:capturing-and-monitoring-the-network:trigger}{trigger packet}'s sequence number (or a hash of it) into a field
of the order it sends, usually its client order identifier, and the capture reads it back.

\begin{method}[Wire-to-wire from a capture]\label{met:nw:capturing-and-monitoring-the-network:w2w}
\begin{enumerate}
\item Tap both market-data lines and the order handoff; timestamp all three against the same clock, on the wire.
\item For each market-data packet keep the time of the first copy seen on either line: the strategy acted on whichever arrived first.
\item Have the strategy carry the trigger's sequence number in every order; for each order, subtract its trigger's time from its own.
\item Report the distribution, not the mean, and break it down against a budget of the path's stages, so that a slow percentile can be
  charged to a stage.
\end{enumerate}
\end{method}

The chapter's tutorial runs this method on a simulated second: 99\,055 market-data packets on two lines, each copy lost independently
with probability 0.2\%, and a strategy that answers 30\% of the packets carrying a trade message with an order, 9\,558 orders in all. The
orders' timing is drawn from \texttt{firm.wirepath}, the extension of One Quant Book~13's tick-to-trade path to the wire
(\cref{fig:nw:capturing-and-monitoring-the-network:budget}). The capture recovers every order's latency exactly by the sequence carried
in the order; the naive rule pairs 47\% of them with the wrong packet.

\begin{omfigure}
\begin{tikzpicture}
\begin{axis}[xbar=0pt,width=11.5cm,height=6.6cm,bar width=5pt,xmode=log,xmin=20,xmax=60000,
  ytick=data,yticklabels from table={figdata/networks/05-capturing-and-monitoring-the-network/budget.csv}{stage},
  y dir=reverse,xlabel={nanoseconds (log scale)},tick label style={font=\small},label style={font=\small},
  xtick={100,1000,10000},xticklabels={100,1\,000,10\,000},
  grid=major,grid style={gray!25},legend columns=2,legend style={font=\footnotesize,at={(0.5,-0.2)},anchor=north,draw=none,
  fill=none,/tikz/every even column/.append style={column sep=8pt}},table/col sep=comma]
\addplot[fill=omDef!55,draw=omDef] table[y expr=\coordindex,x=p50_ns]{figdata/networks/05-capturing-and-monitoring-the-network/budget.csv};
\addplot[fill=omThm!40,draw=omThm] table[y expr=\coordindex,x=p99_ns]{figdata/networks/05-capturing-and-monitoring-the-network/budget.csv};
\legend{median,99th percentile}
\end{axis}
\end{tikzpicture}

\omcaption{The wire-to-wire budget of \texttt{firm.wirepath}: the layer-1 cage of chapter~2 (network in and out), the card's receive and
transmit paths (model values), and the six in-process stages measured by One Quant Book~13 on a laptop (Intel Core Ultra 7 155H, WSL2).
The last bar is their composition, a simulation. Data: \texttt{fig\_capture.py}.}\label{fig:nw:capturing-and-monitoring-the-network:budget}
\end{omfigure}

The budget says where the time goes. At the median the card's two passes (\qty{0.8}{\micro\second} each, model values) and the ring
between the feed handler and the strategy (\qty{0.76}{\micro\second}, measured) are the largest stages; the fibre and devices of the cage
together are \qty{0.56}{\micro\second}. At the 99th percentile the ring owns it, at \qty{32}{\micro\second}: on the laptop that measured it,
a consumer that sleeps or is descheduled wakes late. Composed, the stages give a wire-to-wire median of \qty{4.4}{\micro\second} and a 99th
percentile of \qty{35}{\micro\second}; the capture of the simulated session measures \qty{4.4}{\micro\second} and \qty{38}{\micro\second}, the
difference being the sample.

\begin{omfigure}
\begin{tikzpicture}
\begin{axis}[width=11.5cm,height=5.0cm,xmode=log,xmin=2,xmax=100,ymin=0,ymax=0.16,
  xlabel={wire-to-wire latency from the capture (\unit{\micro\second}, log scale)},ylabel={share of orders},
  ytick={0,0.05,0.1,0.15},yticklabels={0,0.05,0.10,0.15},
  xtick={2,5,10,20,50,100},xticklabels={2,5,10,20,50,100},
  tick label style={font=\small},label style={font=\small},grid=major,grid style={gray!25},table/col sep=comma]
\addplot[ybar,bar width=3pt,fill=omDef!55,draw=omDef] table[x=mid_us,y=share]{figdata/networks/05-capturing-and-monitoring-the-network/w2w_hist.csv};
\end{axis}
\end{tikzpicture}

\omcaption{Simulation: wire-to-wire latency of 9\,558 orders measured from the capture of one simulated second, by the trigger sequence
each order carries (bins of equal width in log scale). The distribution has the long right tail of its slowest stage. Data:
\texttt{fig\_capture.py}.}\label{fig:nw:capturing-and-monitoring-the-network:hist}
\end{omfigure}

\omcode{code/networks/05-capturing-and-monitoring-the-network/python/nw_capture.py}{74}{95}{Reading the capture: each frame's time
from its trailer, the first copy of each market-data packet on either line, and the trigger sequence of each
order.}\label{lst:nw:capturing-and-monitoring-the-network:analyse}

\section{Detecting gaps, drops and microbursts from a capture}

The same capture answers the questions the feed handler of One Quant Book~13 answers from inside, and some it cannot. A gap is a
sequence number missing from a line; a gap on both lines is a message the firm never had, and the capture shows whether it was lost
before the handoff (missing at the tap) or after it (present at the tap, missing in the handler's counters). In the simulated second,
line~A shows 158 gaps and line~B 198, and no message is missing from both: the handler recovered everything from the other line. Bursts
come from the same timestamps: the busiest 100 microseconds of line~A carried \qty{1.24}{\giga\bit\per\second}, nine times its mean of
\qty{0.14}{\giga\bit\per\second}, the same shape as chapter~1's feeds at a smaller scale.

A capture is the firm's evidence as well as its instrument. It shows what the firm actually sent and when, which is what a venue, a
client or a regulator asks about after an incident; European rules ask firms to identify the exact point at which each timestamp is
applied, and a tap on the handoff is the one point that no software change can move.

\section{Monitoring in production: counters, telemetry and alarms}

\begin{definition}[Streaming telemetry]\label{def:nw:capturing-and-monitoring-the-network:telemetry}
\emph{Streaming telemetry}\index{streaming telemetry} is the continuous export by network devices of their counters and state
(per-port traffic, discards, queue occupancy, optical levels, clock state) to a collector, pushed at a chosen interval or on change,
instead of polled by the collector.
\end{definition}

A capture is detailed and expensive: it stores every frame, and a busy cage fills disks in hours. Monitoring runs continuously on
summaries: every device's counters, streamed every second or on change; every server's socket and ring drops; every handler's gap and
staleness counters; every clock's offset. The alarms that matter in a trading network are few and specific:
\begin{itemize}
\item \textbf{discards on an \omterm{def:nw:networking-for-trading:egress}{egress queue}}, which only \omterm{def:nw:networking-for-trading:burst}{microbursts} produce on an underused link (chapter~1);
\item \textbf{gaps on one line} persisting, a line degrading while the other hides it;
\item \textbf{gaps on both lines}, a common cause, and any snapshot recovery;
\item \textbf{a clock out of lock or in \omterm{def:nw:time-synchronisation:holdover}{holdover}}, and offsets approaching the budget (chapter~4);
\item \textbf{wire-to-wire percentiles} drifting from the baseline, computed from the capture every minute.
\end{itemize}
Each alarm names a place on the path, and the capture of the minutes around it is kept for the investigation.

\section{Tutorial: a second of the cage, from the taps}\label{tut:nw:capturing-and-monitoring-the-network:tut}

\begin{tutorial}
\textbf{Goal.} Capture a simulated second at the cage's taps, write it as a pcap file, and measure from it alone what the firm's
programs cannot see. \textbf{End state:} \cref{fig:nw:capturing-and-monitoring-the-network:budget,fig:nw:capturing-and-monitoring-the-network:hist},
the pairing error of the naive rule, and the gap and burst counts quoted in the text.
\begin{tutsteps}
\item \textbf{The budget.} \texttt{nw\_capture.stages()} takes the cage of chapter~2, the card's model stages and Book~13's measured stages;
  \texttt{budget\_rows()} composes them with \texttt{firm.wirepath}.
\item \textbf{The session.} \texttt{session()} builds the feed's frames on both lines with independent losses, the orders with their
  trigger sequence, and appends a trailer to each copy.
\item \textbf{The file.} \texttt{firm.wirecap.write\_pcap(path, records)} writes it (\cref{lst:nw:capturing-and-monitoring-the-network:pcap});
  any tool that reads nanosecond pcap files opens it.
\item \textbf{The analysis.} \texttt{analyse(read\_pcap(path))} pairs orders with their triggers, counts gaps per line and messages lost on
  both, and finds the busiest window (\cref{lst:nw:capturing-and-monitoring-the-network:analyse}); \texttt{python fig\_capture.py} writes the
  charts' data.
\end{tutsteps}
\textbf{What to change next.} Raise the loss to 5\% per line and count the messages lost on both; make the strategy four times faster and
measure how the naive rule's error changes.
\end{tutorial}

\section{Build: capture tools and the wire-to-wire path}\label{bld:nw:capturing-and-monitoring-the-network:wirecap}

\begin{build}
\bfield{Purpose} Two components: \texttt{firm.wirecap}, the firm's capture toolkit, and \texttt{firm.wirepath}, the tick-to-trade path of
One Quant Book~13 extended to the wire, which chapters~7 and~29 use for their latency tables.
\bfield{Interface} \texttt{firm\_wirecap}: \texttt{udp\_frame}, \texttt{tcp\_frame}, \texttt{add\_trailer}, \texttt{strip\_trailer},
\texttt{write\_pcap}, \texttt{read\_pcap}, \texttt{write\_pcapng}, \texttt{read\_pcapng}, \texttt{parse}, \texttt{mold}, \texttt{ORDER},
\texttt{match\_orders}, \texttt{match\_nearest}, \texttt{gaps}, \texttt{bursts}; C++20 reader \texttt{cpp/firm\_wirecap.hpp}.
\texttt{firm\_wirepath}: \texttt{Stage(name, p50, p99, kind)}, \texttt{book13\_software()}, \texttt{wire\_stages(cage, \dots)},
\texttt{card\_stages}, \texttt{hardware\_stage(cycles, mhz)}, \texttt{compose}, \texttt{report}, composing with \texttt{firm.latbudget}.
\bfield{Rules} Timestamps are integers of nanoseconds since midnight UTC; the trailer's time wins over any storage time; the IPv4 checksum is
correct; an order is paired only by the sequence it carries; Book~13's measurements are read, never copied or changed.
\bfield{Acceptance tests} \texttt{code/firm/wirecap/tests/}: frames parse back, checksums hold, trailers round-trip, pcap and pcapng
round-trip, a wrong magic is refused, matching, gaps and bursts on hand-made data, and the shared fixture read identically in Python and
C++. \texttt{code/firm/wirepath/tests/}: the lognormal fit hits its median and 99th percentile, Book~13's stages are read, a cage's path
gives the network stage, and the report's rows.
\bfield{Stretch} Per-port burst detection with the switch's buffer as the threshold; a capture-based order-book rebuild compared, message by
message, with the handler's.
\end{build}

\begin{omsources}
\item \texttt{pcap-savefile(5)} (tcpdump/libpcap); IETF, \emph{PCAP Now Generic (pcapng) Capture File Format},
draft-ietf-opsawg-pcapng-04.
\item Arista, \emph{An Overview of Arista Ethernet Capture Timestamps}; OpenConfig, \emph{gNMI specification}.
\item Commission Delegated Regulation (EU) 2017/574 (RTS~25), article 4.
\end{omsources}

\section{Exercises}

\begin{exercise}[$\star$]\label{exo:nw:capturing-and-monitoring-the-network:1}
A nanosecond pcap record header gives 34\,200 seconds and 1\,500\,000\,321 nanoseconds. Is that a valid record, and what time of day
should the reader report?
\end{exercise}

\begin{exercise}[$\star$]\label{exo:nw:capturing-and-monitoring-the-network:2}
A \omterm{def:nw:capturing-and-monitoring-the-network:trigger}{trigger packet} reaches the tap on line~A at 09:30:00.000\,012\,345 and on line~B \qty{3}{\micro\second} later; the order leaves at
09:30:00.000\,016\,771. What is the wire-to-wire latency? If the copy on line~A had been lost, from which copy would you measure, and
what would change?
\end{exercise}

\begin{exercise}[$\star$]\label{exo:nw:capturing-and-monitoring-the-network:3}
A 10~Gb/s link is tapped in both directions and written to disk with trailers; each frame is stored with a 16-byte record header and
its 16-byte trailer. At a sustained \qty{4}{\giga\bit\per\second} each way with 200-byte frames, how many bytes a second go to disk?
\end{exercise}

\begin{exercise}[$\star\star$]\label{exo:nw:capturing-and-monitoring-the-network:4}
Why does the naive pairing rule fail more often the more frequent the triggers are and the slower the firm is?
\end{exercise}

\begin{exercise}[$\star\star$]\label{exo:nw:capturing-and-monitoring-the-network:5}
Line~A shows 158 gaps and line~B 198 in the tutorial's second, and none is on both. With 0.2\% independent loss per copy and 99\,055
packets, how many packets would you expect to be lost on both lines?
\end{exercise}

\begin{exercise}[$\star\star$]\label{exo:nw:capturing-and-monitoring-the-network:6}
In \cref{fig:nw:capturing-and-monitoring-the-network:budget}, which stage would you attack first to cut the median, and which to cut the 99th
percentile? What would each cost?
\end{exercise}

\begin{exercise}[$\star\star\star$]\label{exo:nw:capturing-and-monitoring-the-network:7}
\emph{Coding.} Make every stage of \texttt{firm.wirepath} four times faster except the network's, run the session again and measure the naive
rule's pairing error. Explain the direction of the change.
\end{exercise}

\begin{exercise}[$\star\star\star$]\label{exo:nw:capturing-and-monitoring-the-network:8}
\emph{Find the flaw.} ``We measure our wire-to-wire latency by timestamping the market-data packet in the handler with the card's hardware
timestamp and the order with the card's transmit timestamp. That is as good as a tap.''
\end{exercise}

\section{Problem: The Two Microseconds Nobody Logged}

\begin{problem}[{Weekend problem --- a latency budget read from the taps}]\label{pb:nw:capturing-and-monitoring-the-network:1}
A firm's strategy logs its in-process time and reports a median of \qty{1.5}{\micro\second}. The capture at the taps (the chapter's
simulated second) reports a wire-to-wire median of \qty{4.4}{\micro\second} and a 99th percentile of \qty{38}{\micro\second}. Use the stage
medians and 99th percentiles of \cref{fig:nw:capturing-and-monitoring-the-network:budget}: network in 263 and out 295~ns; card receive 798
(p99 1\,997) and transmit 801 (p99 2\,000); the six software stages 437, 756, 158, 97, 47 and 33~ns at the median.

\medskip\noindent\textbf{Part I --- Medians.}
\begin{enumerate}
\item What is the sum of the software stages' medians?
\item What is the sum of every stage's median?
\item Why is the sum of medians not the median of the sum?
\item How much of the wire-to-wire median lies outside the program?
\end{enumerate}

\medskip\noindent\textbf{Part II --- The tail.}
\begin{enumerate}[resume]
\item The ring's 99th percentile is \qty{31.9}{\micro\second}. What share of the wire-to-wire 99th percentile is that?
\item Why can a single slow stage own the tail?
\item What did Book~13 find the ring's tail to be on the laptop?
\item What would a production host change?
\end{enumerate}

\medskip\noindent\textbf{Part III --- The capture.}
\begin{enumerate}[resume]
\item Why is the tap's timestamp the right start for the latency, rather than the card's?
\item The naive rule mispaired 47\% of the orders. What would the 99th percentile have looked like with it?
\item What field would you use to carry the trigger in the order?
\item Which alarm would have told the firm that its tail had moved?
\end{enumerate}

\medskip\noindent\textbf{Part IV --- The verdict.}
\begin{enumerate}[resume]
\item State the \emph{named result}: the wire-to-wire median and 99th percentile, the part outside the program, and the stage that owns the
  tail.
\item If the cards' paths were halved, how much would the median fall?
\item If the ring's tail were removed entirely, which stage would own the 99th percentile?
\item What would you need to measure the cards' stages instead of assuming them?
\item Why does the capture belong to the firm's evidence as well as its tools?
\item How much disk does an hour of this capture take, at the tutorial's rates?
\item What would you keep for a month, and what for a day?
\item In one sentence: what does a tap see that a log does not?
\end{enumerate}
\end{problem}

\section{Interview questions}

\begin{interviewq}[$\star$ \iqroles{developer}]\label{iq:nw:capturing-and-monitoring-the-network:1}
How do you measure wire-to-wire latency? What do you need that a program's own logs do not give you?
\end{interviewq}

\begin{interviewq}[$\star\star$ \iqroles{developer}]\label{iq:nw:capturing-and-monitoring-the-network:2}
Where should a capture's timestamp be taken, and why does it matter?
\end{interviewq}

\begin{interviewq}[$\star\star$ \iqroles{developer, researcher}]\label{iq:nw:capturing-and-monitoring-the-network:3}
How do you pair each order with the market-data packet that caused it?
\end{interviewq}

\begin{interviewq}[$\star\star$ \iqroles{developer}]\label{iq:nw:capturing-and-monitoring-the-network:4}
Which network alarms would you set for a trading cage, and why those?
\end{interviewq}

\begin{interviewq}[$\star\star\star$ \iqroles{developer}]\label{iq:nw:capturing-and-monitoring-the-network:5}
Your wire-to-wire 99th percentile doubled overnight; nothing was deployed. How do you find out why?
\end{interviewq}

\begin{interviewq}[$\star\star$ \iqroles{developer, researcher}]\label{iq:nw:capturing-and-monitoring-the-network:6}
What would you use a capture for besides latency?
\end{interviewq}
